\section{Conclusion}
In the present paper, we have derived a formula that relates the short
flow-time behavior of some gauge-invariant local products generated by the
Yang--Mills gradient flow and the correctly-normalized conserved
energy--momentum tensor in the Yang--Mills theory. Our main result is
Eq.~\eqref{eq:(4.29)}. The right-hand side of~Eq.~\eqref{eq:(4.29)} can be
computed by the Wilson flow in lattice gauge theory with appropriate
discretizations of operators, Eqs.~\eqref{eq:(3.5)} and~\eqref{eq:(3.6)} (see,
e.g., Refs.~\cite{Luscher:2010iy,Borsanyi:2012zs}). Here, the continuum
limit~$a\to0$ must be taken first and then the $t\to0$ limit is taken
afterwards; otherwise our basic reasoning does not hold.

Although the formula~\eqref{eq:(4.29)} should be mathematically correct, the
practical usefulness of~Eq.~\eqref{eq:(4.29)} is a separate issue and has to be
carefully examined numerically.\footnote{We hope to return to this problem in
the near future.} Since the lattice spacing~$a$ must be sufficiently smaller
than the square-root of the flow time~$\sqrt{8t}$ for our reasoning to work,
the reliable application of~Eq.~\eqref{eq:(4.29)} will require rather small
lattice spacings. One also worries about contamination by higher-dimensional
operators (i.e., the $O(t)$ terms in~Eqs.~\eqref{eq:(3.7)}
and~\eqref{eq:(3.8)}) and the finite-size effect which we have not taken into
account in the present paper. If our strategy turns to be practically feasible,
it provides a completely new method to compute correlation functions containing
a well-defined energy--momentum tensor. It is clear that the present approach
to the energy--momentum tensor on the lattice is not limited to the pure
Yang--Mills theory although the treatment might be slightly more complicated
with the presence of other fields. The application will then include the
determination of the shear and bulk viscosities (see, e.g.,
Refs.~\cite{Meyer:2007ic,Meyer:2007dy}), the measurement of
thermodynamical quantities (see Ref.~\cite{Giusti:2012yj} and references cited
therein), the mass and the decay constant of the pseudo Nambu--Goldstone boson
associated with the (approximate) dilatation invariance (see
Ref.~\cite{Appelquist:2010gy} and references cited therein), and so on.

It is also clear that our basic idea, that operators defined with lattice
regularization and in the continuum theory can be related through the gradient
flow is not limited to the energy--momentum tensor. For example, it might be
possible to construct an ideal chiral current or an ideal supercurrent on the
lattice, from the small flow-time limit of local products. It would be
interesting to pursue this idea.
